% PROOF REPORT::BATCHFRONTIER1::2026-10-07
\documentclass[11pt]{article}
\usepackage[margin=1in]{geometry}
\usepackage{amsmath,amssymb,amsthm,booktabs,hyperref}
\newtheorem{proposition}{Proposition}
\setlength{\emergencystretch}{3em}
\title{BATCHFRONTIER1: Shared Source Scans for Exact Frontier Merges}
\author{PROOF: concrete WRAPPER realization audit}
\date{7 October 2026}
\begin{document}\maketitle
\paragraph{Status.} COMPUTATIONAL implementation with symbolic accounting supplied for independent review. This run constructs a combined source-backed, shared-setup frontier solver. It improves measured source reads without importing favorable time and space from different implementations. It does not establish a general SAT shrinker or achievable-region contraction. P versus NP remains OPEN.

\section{Recovered components before the build}
The source-backed backends already remove retained residual CNFs, with assignment/trail space and repeated source access charged. FRONTIER-JOINT1 already gives the exact relation for merging satisfying prefix histories: equal assignments on all variables that occur in the suffix imply equal extension possibilities. Its local-rescan backend scans the suffix separately for every current candidate variable.

The fresh STREAMCOMP1 source charges per-access generator time in the recursive inequality $C_i\ge\eta_e+g_e+\theta_e C_j$, and records space floors separately. It is a sufficient conservative simulation, not a lower bound requiring repeated independent work. The fresh ORBITBATCH1 source permits same-parent setup reuse and proves its declared field-address deduplication identity; its field domain $p^r$ is explicitly not original SAT length. The present construction adopts legitimate setup sharing but does not transplant an algebraic field-domain exponent into SAT.

The recursive closure remains conditional on a source-matched shrinking return edge. A clause frontier table is a direct SAT algorithm, not automatically a family of smaller ordinary SAT children. The double wrapper and ARC routes remain available, but this run does not certify such an edge or existential contraction from the bounded-width implementation.

\section{One exact combined realization}
Fix a CNF $F=C_1\land\cdots\land C_m$ in its given clause order, with $L$ literal occurrences, width at most $q$, and maximum boundary width $w$. The boundary after clause $i$ consists of variables occurring both in the prefix and in the suffix. Retain the ordinary viability table on assignments to this boundary.

Choose a positive block size $B$. At the start of a block, form a dictionary whose keys are the incoming boundary variables together with variables occurring in the next at most $B$ clauses. Its size is at most $w+qB$. Scan those clauses once and then the remaining suffix once, recording the last occurrence of each key. Every clause transition in the block now obtains its next boundary by testing recorded last occurrences. No separate suffix scan is performed for each candidate variable or each clause. The dictionary is discarded at the block boundary and constructed afresh for the next block.

The implementation processes the same Boolean table transitions as FRONTIER-JOINT1. It never asks a SAT oracle, calls a reference solver inside production search, reorders the formula, or borrows the propagation backend's resource bounds. A boundary cap returns UNKNOWN, not UNSAT.

\begin{proposition}[Exact boundary and answer preservation]
The batched implementation and the local-rescan implementation produce the same boundaries, viability tables, and answers for every input and block size until a common width-cap exit.
\end{proposition}
\begin{proof}
At a clause inside a block, any candidate next-boundary variable is either an incoming boundary variable or appears earlier in that block. It is therefore in the dictionary. The initial block scan and single suffix scan record its true last occurrence at or after the block start. It belongs to the next boundary precisely when that occurrence is later than the current clause. Thus the boundary test equals local rescanning. Both implementations enumerate the same assignments on the previous boundary plus the current clause variables, check the same old viability entries and clause predicate, and project to the same next boundary. Induction proves equality of every transition and final SAT answer.
\end{proof}

\section{Joint ledger}
There are at most $\lceil m/B\rceil$ blocks. Literal reads obey
\[
R\le \bigl(2+\lceil m/B\rceil\bigr)L.
\]
Two full-input contributions cover block-key collection and actual clause transitions; each shared suffix scan contributes at most $L$. An explicit deterministic balanced dictionary costs logarithmic overhead for each lookup. The Python experiment uses a dictionary and records lookups separately; its measurements do not establish worst-case constant-time hashing.

For $q=O(1)$ and $w=O(\log N)$, a conservative bit/word-model ledger, up to logarithmic factors from identifiers and dictionary operations, is
\[
T=\widetilde O\left(L\lceil m/B\rceil+L+m2^{w+q}\right),
\qquad
S=O\left((w+qB)\log(N+2)+2^w\right)\text{ bits}.
\]
Projection and list search add polynomial factors in $w+q$, hence are polylogarithmic in this regime. Table payload is not total memory. Read-only source access, dictionary storage, all lookups and transitions belong to this same implementation.

For a sparse bounded-width family with $m,L=N^{1+o(1)}$, put $B=N^{\beta+o(1)}$ and $2^w=N^{\omega+o(1)}$, $0\le\beta\le1$. This analysis supplies the envelope
\[
C=\max\{2-\beta,1+\omega\},\qquad
\sigma=\max\{\beta,\omega\}.
\]
It never certifies $C+\sigma<2$, and therefore does not certify the Williams target $C+\sigma<2\cos(\pi/7)$. This observation about the envelope alone is not a lower bound. The next control supplies an actual implementation-specific obstruction.

\section{Hostile chain control}
Let
\[
F_n=\bigwedge_{v=1}^{n-1}(\neg x_v\lor x_{v+1})\land x_1.
\]
It is satisfiable, Horn, and 2-CNF, with frontier width at most two in this order. It never exits early with conflict. For $1\le B\le n$, the implemented full suffix scans cost $\Theta(n^2/B+n)$ literal reads; the dictionary has $\Theta(B)$ entries at its largest block, up to endpoint constants. It follows that this implementation's time-space product is at least $\Omega(n^2)$, already counting only constant cost per source read and per stored entry. With explicit indexed encoding $N=\Theta(n\log n)$, that is $N^{2-o(1)}$.

Thus merely selecting the best block size cannot make this implementation meet the product endpoint on every input. This is not a SAT lower bound: this easy chain has much cheaper specialized algorithms. The repair target is a better source-access plan or semantic preprocessing, with acquisition and storage included. Adding a portfolio may repair this family, but requires a new coverage theorem for all remaining inputs.

\section{Bounded reproducible experiment}
The runner tests 289 formulas at six block sizes: all 256 subsets of full-width signed clauses on three variables; 24 seeded random instances at 6, 8 and 10 variables; five chain sizes through 256; and four empty/contradictory/high-width controls. There are 1728 completed answer checks, six cap exits, and 1734 matched-traversal checks against the prior local-rescan implementation. Every read-count and dictionary-size inequality is also asserted. All checks passed in approximately 0.66 seconds, under an internal 25-second alarm and external 30-second timeout.

\begin{center}\begin{tabular}{rrrr}
\toprule
Block size & Reads, chain 256 & Peak dictionary records & Table bytes\\
\midrule
1 & 66047 & 3 & 8\\
4 & 17087 & 6 & 8\\
16 & 4847 & 18 & 8\\
32 & 2807 & 34 & 8\\
\bottomrule
\end{tabular}\end{center}
The local-rescan baseline reads 195077 literals on that chain; all compared runs perform 2038 transition assignments. At block size 16, the source-read count falls by a factor of about 40.25, with 18 last-occurrence records. This is an operation-count comparison, not a wall-clock speedup or total-memory measurement.

The seed is 20261008. The raw JSON records the exact code SHA256, per-case dimensions, width, table payload, index records, source reads, probes, transitions, answers and cap exits. Code: \texttt{batched\_frontier.py}; data: \texttt{batched\_frontier\_results.json}. The runner imports the prior \texttt{frontier\_wrapper.py} only for baseline comparison and brute-force reference checks.

\section{Corpus, literature, and reconciliation}
Reviewed current GitHub STREAMCOMP1 and ORBITBATCH1 in full; repository-wide search for block/boundary/suffix work; local current source-backed and frontier code/reports; the root synthesis snapshot, governance, CODEX and the repository research skill's experiment mode and guard protocol. The personal-context service returned ``Personal context is unavailable for this conversation,'' so no fresh historical-chat review is claimed. This is targeted corpus reconciliation, not an assertion that every historical document was reread.

Primary literature checked: Allender, Chen, Lou, Papakonstantinou and Tang, \emph{Width-Parametrized SAT: Time--Space Tradeoffs}, Theory of Computing 10 (2014), 297--339, \url{https://www.theoryofcomputing.org/articles/v010a012/}. It studies joint costs on width-parameterized SAT; no unrestricted width-reduction theorem is imported. A current-primary-literature search found no source establishing this project's required general contraction. That limited search is not evidence of a literature-wide impossibility.

Shared block indexing is analogous to batched database lookup and cache amortization. The exact transfer here is only the explicit scan-sharing construction just proved; unrelated database lower bounds are not transferred to SAT. ORBITBATCH1 is adopted for the accounting lesson, not its source-specific field identity. No cryptographic hardness assumption is used. The merge proof is elementary and survives semantic/algebraic oracle replacement; it supplies no nonrelativizing or nonalgebrizing hinge, and no constructive circuit property for a Natural-Proofs claim.

Relative to root MAIN, no final arrow is closed or strengthened. Previously solved boundary exactness and shared-setup accounting are retained. The new executable realization reduces acquisition scans and gives a concrete time-space tradeoff. Independent review of the symbolic bound remains pending. The obsolete statement that every candidate variable needs a separate suffix scan is rejected for this new implementation, not retroactively for the old code. The generic ARC closure, general favorable width/scan guarantee, and source-matched SAT return transformation remain OPEN after matching the actual implementation to the recovered ledgers.

\paragraph{Next target and peer handoff.}
Replace full suffix rescans by a correct event/next-occurrence mechanism with cheaper charged discovery on both the easy chain and adversarial clause orders. RED should audit the implementation-specific product obstruction and cap semantics. LEAD/WRITE should reconcile the new access improvement into main4/main5 without borrowing the fallback's completeness and the frontier's small-space claim simultaneously.
% END_OF_VERIFIED_REPORT::PROOF::BATCHFRONTIER1
\end{document}
